\documentclass{article}
\usepackage[utf8]{inputenc}
\usepackage[T1]{fontenc}
\usepackage{amsmath,amssymb,amsthm,bbm,float,graphicx,geometry,lmodern,mathtools,parskip,setspace,subcaption}
\usepackage[colorlinks]{hyperref}
\usepackage{mathrsfs}
\usepackage{enumerate}
\usepackage{nicefrac}
\usepackage{todonotes}
\usepackage{comment}
\usepackage[notref,notcite]{showkeys}
\newcommand{\e}{\mathrm e}

%%%Tikz
\usepackage{tikz}
\usetikzlibrary{backgrounds}
\usetikzlibrary{patterns}
\usetikzlibrary{positioning, shapes.geometric}

%%% Captionsize
\usepackage{caption}
\captionsetup{font=footnotesize}
\usepackage{graphicx}
%%%%%%Formatting%%%%%%
\graphicspath{{Images/}}%put any images in this file
\parindent0pt%paragraph indent
\allowdisplaybreaks%breaks long equations onto multiple pages is needed
%%%%%%Commands%%%%%%
\DeclarePairedDelimiter{\set}{\lbrace}{\rbrace}
\newcommand{\N}{\mathbb{N}}
\DeclareMathOperator{\argmax}{arg \ max}
%%%%%%Environments%%%%%%

\newtheorem{theorem}{Theorem}[section]
\newtheorem{lemma}[theorem]{Lemma}
\newtheorem{corollary}[theorem]{Corollary}
\newtheorem{prop}[theorem]{Proposition}
\newtheorem{conj}[theorem]{Conjecture}

\theoremstyle{definition}
\newtheorem{definition}[theorem]{Definition}
\newtheorem{remark}[theorem]{Remark}

%%%%%%%% symbols and short cuts%%%%%%%
\renewcommand{\P}{\mathbb{P}}
\newcommand{\x}{\mathbf{x}}
\newcommand{\y}{\mathbf{y}}
\newcommand{\1}{\mathbbm{1}}
\newcommand{\G}{\mathscr{G}}
%\renewcommand{\H}{\mathscr{H}}
\newcommand{\C}{\mathscr{C}}
\newcommand{\X}{\mathbf{X}}
\newcommand{\E}{\mathbb{E}}
\newcommand{\0}{\mathbf{o}}
\newcommand{\R}{\mathbb{R}}
\renewcommand{\d}{\mathrm{d}}
\newcommand{\z}{\mathbf{z}}
\newcommand{\w}{\mathbf{w}}
\renewcommand{\v}{\mathbf{v}}
\renewcommand{\u}{\mathbf{u}}
\newcommand{\Z}{\mathbb{Z}}
\newcommand{\T}{\mathcal{T}}
\newcommand{\U}{\mathcal{U}}
\newcommand{\cX}{\mathcal{X}}
\newcommand{\cG}{\mathcal{G}}


%Commands Bas
\newcommand*{\be}{\begin{equation}}
\newcommand*{\ee}{\end{equation}}
\newcommand*{\ba}{\begin{aligned}}
\newcommand*{\ea}{\end{aligned}}

\newcommand{\D}{\mathbb D}
\newcommand{\cB}{\mathcal B}
\newcommand{\cP}{\mathcal P}
\newcommand{\cU}{\mathcal U}
%\newcommand{\ind}{\mathbbm 1}
\newcommand{\wt}{\widetilde}
\newcommand{\cR}{\mathcal R}
\newcommand{\sm}{\mathrm{sm}}
\newcommand{\basc}[1]{{\color{blue}{ \bf [~Bas:\ }\emph{#1}\textbf{~]}}}

%Packages Bas
%\usepackage{enumitem}
\usepackage{bbm}

%%%%%%Bibliography%%%%%%
\usepackage[style = numeric, sorting=nyt, url = false, abbreviate=false, maxbibnames=9, sortcites=true, doi = true, backend = biber, giveninits = true]{biblatex}
\renewbibmacro{in:}{\ifentrytype{article}{}{\printtext{\bibstring{in}\intitlepunct}}}
\bibliography{bib.bib}

\title{Percolation in lattice $k$-neighbor graphs}
\author{Benedikt Jahnel, Jonas K\"oppl, Bas Lodewijks, Andr\'as T\'obi\'as}
\date{\today}

\begin{document}

\maketitle

\begin{abstract}TTT
\end{abstract}

\section{Introduction}


\section{Setting and main results}\label{sec:set}

Consider the $d$-dimensional cubic lattice $\mathbb Z^d$. We are interested in the percolation behavior of the {\em $k$-neighbor graph} ($k$-nG) in which each vertex independently chooses uniformly $k\le 2d$ of its $2d$ nearest neighbors. We distinguish three different types of $k$-nGs:
\begin{enumerate}
\item The {\em directed $k$-nG} ($k$-DnG) in which we open a directed edge from the vertex to each of the $k$ chosen neighbors. 
\item The {\em undirected $k$-nG} ($k$-UnG) in which we open an undirected edge from the vertex to each of the $k$ chosen neighbors.
\item The {\em bidirectional $k$-nG} ($k$-BnG) in which we open an undirected edge between two vertices whenever both choose the other one as one of its $k$ chosen neighbors.
%\item The {\em exclusively unidirectional $k$-nG} or ($k$-XnG) in which we open an undirected edge between two vertices whenever exactly one of them chooses the other one as one of its $k$ chosen neighbors. The letter X refers to the ``XOR'' (exlusive ``or'') in the edge-drawing rule.
\end{enumerate}

\color{blue}
\begin{remark}
These graph satisfy the following complementary relations.
\begin{itemize}
    \item For $1 \le k \leq 2d$, the directed graph with vertex set $\Z^d$ having a directed edge from $u \in \Z^d$ to $v \in \Z^d$ if and only if $u$ and $v$ are nearest neighbors and $(u,v)$ is not an edge in the $k$-DnG is a $2d-k$-DnG.
    \item For $1 \leq k \leq 2d$, the graph-theoretic complement of the $k$-UnG is a $2d-k$-BnG. (Hence, the same is also true with the words ,,UnG'' and ,,BnG'' swapped.)
    %\item For $1 \leq k \leq 2d$, the graph-theoretic complement of the $k$-XnG is the disjoint union of two correlated graphs, their marginal distributions being the one of a $k$-BnG and a $2d-k$-BnG. 
\end{itemize}
\end{remark}
\color{black}
We are interested in the behavior of the percolation probabilities 
$$
\theta^{\rm D}(k,d)=\mathbb P(o\rightsquigarrow\infty \text{ in the $k$-DnG on $\mathbb Z^d$}),
$$
respectively $\theta^{\rm U}(k,d)$ and $\theta^{\rm B}(k,d)$, %and $\theta^{\rm X}(k,d)$, 
where $o\rightsquigarrow\infty$ represents the event that there exists a path along open edges from the origin to infinity. Let us note that the percolation probabilities are clearly non-decreasing functions in $k$. 

Let us mention that the outdegree in $k$-DnG is $k$ almost surely. However, the degree in $k$-BnG is a binomial random variable with parameters $k$ and $k/2d$, i.e., the expected degree is given by $k^2/2d$ which is less than $k$, unless $k=2d$. In case of the $k$-UnG, the degree is distributed according to $k+D$, where $D$ is a binomial random variable with parameters $2d-k$ and $k/2d$, and hence, the expected degree is given by $k(4d-k)/(2d)$.

\subsection{Results for the directed $k$-neighbor graph}
\begin{lemma}\label{lem:surenoperc}
We have for all $d\ge 1$ that $\theta^{\rm D}(1,d)=0$.
\end{lemma}
\begin{proof}
Indeed, if there exists a path starting from the origin and reaching an $\ell_1$-distance $n$ from the origin, then this path is unique and has at least $n$ steps. But such a path exists with probability at most $(2d)^{-n}$, which converges to zero as $n$ tends to infinity.
\end{proof}
%Note that the $k$-nnG has a natural dual in which each existing directed edge is removed and all missing edges are present. This graph is edge-disjoint but not vertex disjoint and it is in distribution equivalent to the $(2d-k)$-nnG. 
% \begin{lemma}\label{lemma2}
% We have that $\theta(2d-1,d)>0$.
% \end{lemma}
% \begin{proof}
% Note that by ergodicity, since $\theta(1,d)=0$ in the $(2d-1)$-nnG in $d$ dimensions, there exist only finite components in the $1$-nnG in $d$ dimensions. But this means that the $(2d-1)$-nnG in $d$ dimensions percolates and thus $\theta(2d-1,d)>0$. 
% \end{proof}
\begin{lemma}\label{lem:sureperc}
    We have that $\theta^{\rm D}(k,d)=1$ for all $k\ge d+1$.
\end{lemma}
In particular, from Lemma~\ref{lem:surenoperc} and Lemma~\ref{lem:sureperc} we see that for $d=2$ the critical $k$ is given by $2$. Simulations suggest that $\theta^{\rm D}(2,2)>0$. We consider a proof of this statement a major success. 
\begin{proof}
    We will use a growth argument and show that the maximal distance to the origin is strictly increasing between generations. To make this precise, denote by $G_n$ the new vertices discovered in the $n$-th step, where we start with $G_0 = \{0\}$ and at every step $n \to n+1$, each vertex in $G_n$ chooses $k$ of its neighbours uniformly at random as potential successors and $G_{n+1}$ is then the set of all potential successors that have not been discovered before for any $m  \leq n$. For $x\in\Z^d$, we let $\|x\|_1$ denote the $\ell_1$-distance between $x$ and the origin. We show by induction that 
    \begin{align}\label{eqn:growth-property}
         \forall n  \in \N:   \quad \max_{x \in G_n}\|x\|_1 < \max_{x \in G_{n+1}}\|x\|_1.  
    \end{align}
    This then implies that there exists an infinitely long directed path starting at the origin. 
    
    For $n=0$ the inequality in~\eqref{eqn:growth-property} is clear. For the induction step, let $x \in G_n$ be a vertex (possibly non unique) that achieves the maximum. Then $x$ has $d$ neighbours $y$ with $x_1 < y_1$. If $k  \geq 2d-d+1$, then we necessarily have to choose one of these $y$ as a potential vertex for the new generation. This $y$ cannot have been in any of the previous generations by the induction hypothesis. Therefore we have $y \in G_{n+1}$ and~\eqref{eqn:growth-property} follows. 
    \end{proof}

Let us turn our attention towards the supercritical percolation phase. We have the following result. For this, consider two independent random walks $S$ and $S'$ with $S_0=e_1$ and $S'_0=e_2$, where $(e_1,\dots, e_d)$ denote the unit vectors. The walks move as simple random walks in the positive direction, i.e., $S_{n+1}=S_n+e_i$ where $i$ is uniformly chosen from $\{1,\dots ,d\}$. We define by
\begin{align*}
\rho(d)=\P(\exists m\ge 1\colon S_m=S'_m)\in(0,1]
\end{align*}
the meeting probability of the walkers. 
\begin{theorem}
Let $k\le d$ be given, then, whenever $\rho(d)<(2d-1)(k-2)/(2(d-1)(k-1))$, we have that $\theta^{\rm D}(k,d)>0$.
\end{theorem}
From this we see that for $k=2$ there is no criterion. For $k=3$ we have $\rho(d)<(2d-1)/(4(d-1))$, and for $d=3$ we have $\rho(d)<5/8$, for $d=4$ we have $\rho(d)<7/12$. For some preliminary calculations for $\rho(d)$, see below. 
\begin{proof}
We follow Cox--Durrett. Let $\mathcal R_n$ denote the set of directed paths from the origin to level $n$ (that is, all vertices at $\ell_1$ distance $n$), and let $N_n$ be the number of open paths in $\mathcal R_n$. \basc{Is any path from the origin to level $n$ in $\mathcal R_n$, or does it need to be a path with strictly increasing $\ell_1$ distance?} Then $W_n=(2/k)^nN_n$ is a martingale with respect to $\mathcal F_n$, the sigma-algebra generated by the first $n$ steps of a simple random walk $T$ with $T_0=o$ and $T_{n+1}=T_n+e_i$ where $i$ is uniformly chosen from $\{1,\dots ,d\}$. Indeed, using the fact that the expected number of outgoing edges \basc{ we should make clear here what we mean with `outgoing edges'}  is given by the expectation of a hypergeometric random variable, i.e., 
\begin{align*}
\sum_{l=0}^k l\frac{\binom{d}{l}\binom{d}{k-l}}{\binom{2d}{k}}=\frac{k}{2},
\end{align*}

\begin{comment} 
\basc{Observe that this is only true for vertices $v$ such that the number of neighbours with larger (resp.\ smaller) $\ell_1$-distance equals $d$ (resp.\ $d$). There are, however, vertices for which this is not the case. In any dimension $d\geq 2$ and any $n\in\N$, there are always $2d$ vertices at distance $n$ such that they have $2d-1>d$ many neighbours at $\ell_1$-distance $n+1$ and only $1$ neighbour at $\ell_1$-distance $n-1$. In particular, when $d=k=2$,
\be 
N_1=k \text{ almost surely, and } \E[N_2]=12\Big(\frac{k}{2d}\Big)^2=3\neq 2. 
\ee
Similarly, 
\be \ba
\E[N_3]&=4\cdot 3\Big(\frac{k}{2d}\Big)^3+4\cdot 2\cdot 2\Big(\frac{k}{2d}\Big)^3=28\Big(\frac{k}{2d}\Big)^3=\frac72\neq 2,\\
\E[N_4]&=4\cdot 3\Big(\frac{k}{2d}\Big)^4 +8\cdot 3\cdot 2 \Big(\frac{k}{2d}\Big)^4=\frac{15}{4}\neq 2. {\color{red} \text{this may be incorrect}}
\ea \ee  
As $n$ grows, $\E[N_n]$ will become closer to $2$ (in the case $k=2$, or to $k(k/2)^n$ in general, see also my comment further on) I think, as the contribution of the paths from the $4$ vertices with $3>d$ many neighbours decreases relatively compared to the contribution of the other paths. \\
In general (with $d=2$ and $k\leq d$), at $\ell_1$-distance $n$ there are $4(n-1)$ vertices who have, in expectation, $k/2$ many neighbours that increase their $\ell_1$-distance w.r.t.\ their parent, and $4$ vertices with, in expectation, $3k/4$ many neighbours that increase their $\ell_1$-distance w.r.t.\ their parent.\\
As a result, $W_n$ is a sub-martingale rather than a martingale. 
}
\end{comment}
we can use independence to see that 
\begin{align*}
\E[N_{n+1}|\mathcal F_n]=\frac{k}{2}N_n.
\end{align*}
Since $\E[W_1]=1$, we can apply the martingale convergence theorem, to ensure the existence of a random variable $W$ with 
\begin{align*}
W_n\to W\qquad \text{almost surely}.
\end{align*}
Using the second-moment method it now suffices to show that $\limsup_{n\uparrow\infty}\E[W_n^2]<\infty$.

For this, note that 
\begin{align*}
\E[N_n^2]=\sum_{s,t\in \mathcal R_n}\P(s,t\text{ open})=\sum_{s,t\in \mathcal R_n}p^{K(s,t)}q^{L(s,t)}p^{2(n-K(s,t)-L(s,t))}
\end{align*}
where $p=k/2d$ denotes the probability that a given edge is open, $q=k(k-1)/(2d(2d-1))$ denotes the probability that two different edges that emerge from the same vertex are open, $K(s,t)$ is the number joint edges in two paths $s$ and $t$, and $L(s,t)$ is the number of vertices $x$ in $s$ and $t$ such that $s$ and $t$ do not have a joint edges after $x$. \color{green} ** \color{black} This implies that 
\begin{align*}
\lim_{n\uparrow\infty}\E[W_n^2]=\E'\Big[p^{-K}\Big(\frac{q}{p^2}\Big)^L\Big]=\E[W_n^2]=\E'[a^Kb^L]
\end{align*}
\basc{The third equality should be omitted, right?}
where $\E'$ is the expectation associated to two independent simple random walks $T$ and $T'$ with random $K,L$ as defined above, and $a=2d/k$ and $b=2d(k-1)/(k(2d-1))$. Now, using the definition of $\rho(d)$ from above, we have that 
\begin{align*}
\E'[a^Kb^L]&=\sum_{\ell\ge 0}\sum_{n\ge 0}a^\ell b^n\P'(K=\ell,L=n)\\
&=(1-\rho(d))\sum_{\ell\ge 0}\sum_{n\ge 0}\rho(d)^n\Big(\frac{1}{d}\Big)^\ell\Big(1-\frac{1}{d}\Big)^n\binom{n+\ell}{n}a^\ell b^n\\
&=(1-\rho(d))\sum_{\ell\ge 0}\sum_{n\ge 0}\Big(\frac{a}{d}\Big)^\ell\Big(\frac{\rho(d)b(d-1)}{d}\Big)^n\binom{n+\ell}{n}\\
&=(1-\rho(d))\sum_{\ell\ge 0}\sum_{n\ge 0}\Big(\frac{a}{d}\Big)^{n+\ell}\Big(\frac{\rho(d)b(d-1)}{a}\Big)^n\binom{n+\ell}{n}\\
&=(1-\rho(d))\sum_{n\ge 0}\sum_{\ell\ge n}\Big(\frac{a}{d}\Big)^{\ell}\Big(\frac{\rho(d)b(d-1)}{a}\Big)^n\binom{\ell}{n}\\
&=(1-\rho(d))\sum_{\ell\ge 0}\sum_{n=0}^\ell\Big(\frac{a}{d}\Big)^{\ell}\Big(\frac{\rho(d)b(d-1)}{a}\Big)^n\binom{\ell}{n}\\
&=(1-\rho(d))\sum_{\ell\ge 0}\Big(\frac{a}{d}\Big)^{\ell}\sum_{n=0}^\ell\Big(\frac{\rho(d)b(d-1)}{a}\Big)^n\binom{\ell}{n}\\
&=(1-\rho(d))\sum_{\ell\ge 0}\Big(\frac{a}{d}\Big)^{\ell}\Big(\frac{\rho(d)b(d-1)}{a}+1\Big)^\ell.
\end{align*}
Hence, the criterion is given by 
\begin{align*}
1&>\frac{\rho(d)b(d-1)+a}{d}
\end{align*}
which is equivalent to
\begin{align*}
\rho(d)<\frac{d-a}{b(d-1)}=\frac{d-2d/k}{(d-1)2d(k-1)/(k(2d-1))}=\frac{(2d-1)(k-2)}{2(d-1)(k-1)}
\end{align*}
\end{proof}

% Positive results about percolation in lower dimensions ``propagate'' to higher dimensions thanks to the following lemma. IS THIS REALLY GIVING NEW INFORMATION?
% \begin{lemma}\label{lemma-plane}
% We have that $\theta^{\rm D}(2d'-k,d')\geq \theta^{\rm D}(2d-k,d)$ for all $k\in \{ 1,2,3,d-1 \}$ and $d' \geq d \geq 1$. 
% \end{lemma}
% \begin{proof}
% Let us consider the graph whose vertex set is the intersection of $\Z^{d'}$ with a $d'$-dimensional coordinate plane and its edge set is the set of the edges of the $(2d'-k)$-nnG pointing from one point of the coordinate plane to another. In this graph, every vertex has outdegree at least $2d-k$, and hence the graph stochastically dominates a $2d-k$-nnG in $d$ dimensions, which implies the lemma.
% \end{proof}

%Note that, applying Lemma~\ref{lemma-plane} to $d=2$ and $k=1$, we obtain $\theta(2d'-1,d')\geq \theta(3,2)$, and thus $\theta^{\rm D}(2,2)>0$ follows as soon as one verifies it for $d=2$. Further, Lemma~\ref{lemma-plane} applied to $d=k=2$ implies that if Theorem~\ref{thm1} is true, then $\theta(2d'-2,d')>0$ holds for all $d'\geq 2$. 

\color{green} Alternative statement, whose proof is initially the same as the proof of Theorem 2.3 starting from the ** above. \color{black}

\begin{theorem}\label{thm-CoxDurrett}
$\theta^{\rm D}(k,d)>0$ holds if $d \geq 2$ and $k \geq 4$, or if $d \geq 5$ and $k=3$. 
\end{theorem}

\begin{proof} \color{green} ** \color{black} Now note that $q<p^2$, which implies that
\[ \E[N_n^2] \leq \sum_{s,t\in \mathcal R_n} p^{2n-K(s,t)}. \]
The right-hand side is precisely the same as the right-hand side of the display below Equation~(2$\cdot$7) in~\color{red}Cox--Durrett\color{black}. Thus, a verbatim application of the arguments of \color{red} Section 2 of Cox--Durrett \color{black} implies that 
\[ \lim_{n\to\infty} \E W_n^2 <\infty \]
holds if and only if 
\[ p = \frac{k}{2d} < \varrho(d), \]
where for two independent (simple, symmetric, nearest neighbor) random walks $\widetilde S=(\widetilde S_n)_{n\in\N_0}$ and $\widetilde S'=(\widetilde S'_n)_{n\in\N_0}$ on $\Z^d$ started from the origin, we define
\[ \varrho(d) = \P(\exists m\geq 0 \colon \widetilde S_m = \widetilde S'_m \text{ and } \widetilde S_{m+1} = \widetilde S'_{m+1}). \]
Let further $\tau_d=\inf \{ m \geq 0 \colon \widetilde S'_m \text{ and } \widetilde S_{m+1} = \widetilde S'_{m+1} \}$.
According to \color{red} Cox--Durrett, \color{black} we have for all $d$ that
\begin{equation}\label{firstfew} \P(\tau_d=0)=\frac{1}{d}, \qquad \P(\tau_d=1)=0, \qquad \P(\tau_d=2)=d^{-3}-d^{-4}, \end{equation}
and for $3 \leq k \leq d$,
\[ \P(\tau=k) \leq d^{-k} k!, \]
while for $j\geq 1$ and $k >jd$,
\begin{equation}\label{bigk} \P(\tau=k) \leq d^{-1} (2\pi d)^{1/2} \Big( \frac{\e^{-1/13}}{\sqrt{2\pi}} \Big)^d j^{-\frac{1-d}{2}}, \end{equation}
so that
\begin{equation}\label{CDUB} \begin{aligned} \varrho(d)=\P(\tau_d<\infty) &\leq \frac{1}{d} + \frac{1}{d^3} -\frac{1}{d^4} + \sum_{k=3}^d d^{-k} k! +  (2\pi d)^{1/2} \Big( \frac{\e^{-1/13}}{\sqrt{2\pi}} \Big)^d \sum_{j=1}^{\infty} j^{-\frac{1-d}{2}} \\ & = \frac{1}{d} + \frac{1}{d^3} -\frac{1}{d^4} + \sum_{k=3}^d d^{-k} k! +  (2\pi d)^{1/2} \Big( \frac{\e^{-1/13}}{\sqrt{2\pi}} \Big)^d \zeta\Big( \frac{d-1}{2} \Big), \end{aligned} \end{equation}
where $\zeta$ denotes the Riemann zeta function. (Note that since $\varrho(d)>1/d$, this approach cannot give any positive percolation result for $k=2$.)

%Therefore, it suffices to find the smallest $d$ such that the upper bound of~\eqref{CDUB} is less than $3/(2d)$, and this will imply that $\theta^{\rm D}(k',d') >0$ for all $k' \geq k$ and $d' \geq 3$.
%Note that this will not give any positive percolation result for $k=2$ because $k/2d = 1/d$ is smaller than the right-hand side of~\eqref{CDUB}. The next table shows the values of the right-hand side for $d=4,5,6,7$ and the respective smallest values of $k$ such that these bounds imply $\theta^{\rm D}(k,d)>0$.

The next table shows the smallest values of $k$ such that the right-hand side of~\eqref{CDUB} (which we computed numerically) is less than $k/(2d)$, for $d=4,5,6,7$.
\begin{center}
    \begin{tabular}{|c|c|c|c|c|} \hline
       &  $d=4$ & $d=5$ & $d=6$ & $d=7$ \\ \hline
        Value of the r.h.s.\ of~\eqref{CDUB}   & 0.693093  & 0.394622  & 0.268615 & 0.199707  \\ \hline
         Smallest $k$ such that r.h.s.\ of~\eqref{CDUB} is less than $\tfrac{k}{2d} $ & 6 & 4 & 4 & 3 \\ \hline 
    \end{tabular}
\end{center}
Via additional numerical computations we see that the right-hand side of~\eqref{CDUB} minus $k/(2d)$ is negative for any $k\geq 3$ if $d=7,8,9,10$ and for any $k\geq 4$ if $d=5,\ldots,10$. Given this, it is elementary to check that it stays negative for any $k \geq 3$ if $d \geq 10$.  \color{red} OK? For this it suffices to show that the rhs of~\eqref{CDUB} for dimension $d+1$ is at most $d/(d+1)$ times the same for dimension $d$, which I think is possible if $d \geq 10$. We just have to compensate the additional term $(d+1)^{-(d+1)} (d+1)!$ in the sum in dimension $d+1$. \color{black} 

Hence, it follows that $\theta^{\rm D}(5,d) >0$ for all $d \geq 5$. Moreover, Lemma~\ref{lem:sureperc} implies that $\theta^{\rm D}(5,d)=1$ holds for $d=3,4$. Thus, we have $\theta^{\rm D}(k,d) >0$ whenever $k\geq 5$ (given that it is well-defined, i.e., $k\leq 2d$). The table also tells that  $\theta^{\rm D}(4,d)>0$ for all $d \geq 5$. Moreover, $\theta^{\rm D}(4,2)=1$ is clear and $\theta^{\rm D}(4,3)=1$ follows from Lemma~\ref{lem:sureperc}. Finally, $\theta^{\rm D}(3,d)>0$ for $d \geq 7$ can be seen in the table (and $\theta^{\rm D}(3,2)=1$ is another consequence of Lemma~\ref{lem:sureperc}). Therefore, in order to finish the proof of the theorem, we have to treat three special cases separately: $(k,d)=(4,4)$, $(k,d)=(3,5)$ and $(k,d)=(3,6)$. 

As for the first one, we compute (using a computer) the precise values of $\P(\tau_4=k)$ for $k=3,4,5$ via generating all length-6 random walk trajectories and computing how many of them belong to these events respectively (for $\P(\tau_4=l)$, $l=0,1,2$, we of course obtain the same precise values as \color{red} Cox and Durrett \color{black}). This yields
\[ \P(\tau_4=3)=\frac{3}{512}, \qquad \P(\tau_4=4)=\frac{279}{65\ 536}, \qquad \P(\tau_4=5)=\frac{831}{262\ 144}. \]
Thus, using~\eqref{firstfew} and~\eqref{bigk},
\[ \begin{aligned} \P(\tau_4<\infty) & \leq \frac{1}{4} + \frac{3}{256} + \frac{3}{512} + \frac{279}{65\ 536} + \frac{831}{262\ 144} + \frac{3}{4} \sqrt{8\pi} \Big( \frac{\e^{-1/13}}{\sqrt{2\pi}} \Big)^4 + \sum_{j=2}^{\infty} \sqrt{8\pi} \Big( \frac{\e^{-1/13}}{\sqrt{2\pi}} \Big)^4 \sum_{j=1}^{\infty} j^{-\frac{3}{2}} \\ & \approx 0.495542 < 0.5 = \frac{k}{2d}. \end{aligned}\]
It follows that $\theta^{\rm D}(4,4)>0$.

With similar methods we can show that $\theta^{\rm D}(3,5)>0$ and $\theta^{\rm D}(3,6)>0$. Indeed, for $d=5$ we obtain via explicit computer computation
\[ \P(\tau_5=3)=\frac{44}{15\ 625}, \qquad \P(\tau_5=4)=\frac{712}{390\  625}, \qquad \P(\tau_5=5)=\frac{12\ 136}{9\ 765\ 625}. \]
This already implies that  $\theta^{\rm D}(3,5)>0$ because
\[ \P(\tau_5<\infty) \leq \frac{1}{5} + \frac{4}{625} + \frac{44}{15\ 625} + \frac{712}{390\  625} + \frac{12\ 136}{9\ 765\ 625} + \sqrt{10\pi} \Big( \frac{\e^{-1/13}}{\sqrt{2\pi}} \Big)^5 \zeta(2) \approx 0.275703, \]
\color{red} (for fun we could write that the last summand is $\sqrt 5/(24 \e^{5/13})$) \color{black} which is less than $k/(2d)=0.3$. For $d=6$, determining the value of $\P(\tau_6=k)$, $k=3,\ldots,6=d$ numerically would be computationally quite demanding. However, we see from the table that the right-hand side of~\eqref{CDUB} is just slightly larger than $k/(2d)=1/4=0.25$. It turns out that in this case, it suffices to compute $\P(\tau_6=3)$ numerically, which equals $\frac{35}{23328}$. Estimating $\P(\tau_6>3)$ according to~\eqref{CDUB}, we obtain
\[ \P(\tau_6<\infty) \leq \frac{1}{6} + \frac{5}{1296} + \frac{35}{23\ 328} + \sum_{k=4}^6 6^{-k} k! + \sqrt{12\pi} \Big( \frac{\e^{-1/13}}{\sqrt{2\pi}} \Big)^6 \zeta(5/2) \approx 0.242338 < 0.25=\frac{k}{2d}. \]
\end{proof}

\color{green} After this, the remaining unresolved cases are: $k=2$ and $d \geq 2$; $k=3$ and $d=3,4$. For $(k,d)=(3,4)$ it is not impossible that some Cox--Durrett style argument works, but so far I only see a solution via computing $\P(\tau_4=k)$ until ca.\ $k=12$, which is too much for our numerical capacity. For $d=3$ and for $k=2$ the Cox--Durrett (or Bene's similar) argument doesn't work. \color{black}


Let us collect some observations for $\rho(d)$. As in Cox--Durrett, we can define 
$$
\tau_d=\inf\{n\ge 1\colon S_n=S'_n\}
$$
and it suffices to evaluate
$$
\rho(d)=\sum_{n\ge 1}\P(\tau_d=n),
$$
where $\P(\tau_3=1)=1/9$ and $\P(\tau_3=2)=5/81$. On a general level, using the Cox--Durrett approach, we have that 
\begin{align*}
\P(\tau=1)&=d^{-2}\\
\P(\tau=n)&\le \Big(1-\Big(\frac{(d-1)}{d}\Big)^n\Big)d^{-n}n!,\qquad\text{ for } 2\le n\le d \\
\P(\tau=n)&\le d^{-n}\frac{n!}{\big(\lfloor n/d\rfloor!\big)^d},\qquad\text{ for } d+1\le n.
\end{align*}
Here we used that, in order for the random walks to meet in one step, they have to choose $e_1$ and $e_2$ correspondingly. For a meeting with $n\le d$ steps, we note that if the walker $S$ always walks away from the walker $S'$, then, $(d-1)^n$ paths of $S$ can not be met by $S'$, which leads to the first factor. The second factor is from Cox--Durrett. Also the last line is from Cox--Durrett. 


\bigskip
It would be desirable to proof the following conjecture.
\begin{conj}
For all $k\ge 1$ we have that, whenever $\theta^{\rm D}(k,d)>0$, then also $\theta^{\rm D}(k,d+1)>0$. 
\end{conj}
If this is true, then $k=2$ would always be critical. It is not to be expected that $d\mapsto\theta^{\rm D}(k,d)$ is non-decreasing since $\theta^{\rm D}(3,2)=1$ but heuristically $\theta^{\rm D}(3,3)<1$. IDEA: USE COUPLING. IS $n\mapsto \theta^{\rm D}(k+n,d+n)$ NON-DECREASING? 

It would be desirable to proof the following conjecture.
\begin{conj}
$\theta^{\rm D}(2,2)>0$. 
\end{conj}
IDEA: USE PEIERLS ARGUMENT

\subsection{Results for the undirected $k$-neighbor graph}
We need some results here.
\begin{conj}
We have that $\theta^{\rm U}(k,d)>0$ for all $k,d\ge 2$.  
\end{conj}

\medskip

It is clear that $\theta^{\rm U}(k,d) \geq \theta^{\rm D}(k,d)$ for any $k,d$ because any directed edge in the $k$-DnG corresponds to an undirected edge in the $k$-UnG (and at most two directed edges can correspond to the same undirected edge). Additionally, $k=1$ is not sufficient for percolation even in the undirected sense.

\begin{prop}
    We have that $\theta^{\rm U}(1,d)=0$ for all $d \geq 1$.
\end{prop}
\begin{proof}
Assume for a contradiction that the assertion of the proposition is wrong, and let $A$ then denote the event of positive probability that there exists an infinite path starting from $o$ consisting of edges of the $1$-UnG. %First, we claim that if such a path exists, it can contain at most one bidirectional edge. Indeed, two bidirectional edges cannot be adjacent to each other because this would mean that the vertex separating them has outdegree at least 2 in the $1$-DnG. Also, vertices belonging to the bidirectional edge have no outgoing edge but the one contributing to the bidirectional edge. This way, if $Y_0=o, Y_1,Y_2,\ldots$ is the enumeration of the vertices among an infinite path and for some $i \geq 0$, $(Y_i,Y_{i+1})$ is a bidirectional edge, then either $Y_{i+2}=Y_i$ and $Y_i$ has an additional incoming edge in the $1$-DnG coming from $Y_{i+3} \neq Y_{i+1}$, or $Y_{i+2} \neq Y_i$ and $Y_{i+1}$ has an additional incoming edge coming from $Y_{i+2}$. Now, 
%such a path, then if there is a bidirectional edge, say $(Y_{i},Y_{i+1})$, then the only way that this edge is in the infinite path is that $Y_i$ or $Y_{i+1}$ has an additional incoming edge from some other vertex
%there also exists one with the same property that does not contain any bidirectional edge apart from the very first one (which is then adjacent to $o$). Indeed, in any infinite path of the $1$-UnG, there cannot be two bidirectional edges adjacent to each other because this would mean that the vertex between them has outdegree at least 2. Thus, taking any infinite path starting from $o$, this can only contain isolated bidirectional edges (i.e., ones not adjacent to any other one). Hence, if $Y_0=o,Y_1,Y_2,\ldots$ is an infinite path of the $1$-UnG and for some $k \geq 0$, $(Y_k,Y_{k+1})$ is a bidirectional edge, it must be the case that either $k=0$, or $(Y_{k-1},Y_k)$ and $(Y_{k+2},Y_{k+1})$ is an edge of the $1$-DnG. In either of the cases, if we leave out the bidirectional edge $(Y_k,Y_{k+1})$ from the infinite path, the path still remains infinite, and this can be done simultaneously with all bidirectional edges of the path. This implies the claim.
let $X_0=o,X_1,X_2,\ldots$ be any infinite path of the $1$-UnG starting from $o$.
%that contains no bidirectional edges apart from $(X_0,X_1)$ and is defined (on the event $A$) as follows. First, consider the infinite cluster of $o$ in the $1$-UnG and remove all bidirectional edges of it that are not adjacent to $o$, along with the vertices that are only adjacent to a bidirectional edge of the cluster. Now, we define $X_1,X_2,\ldots$ recursively as follows. If $o$ is not adjacent to a bidirectional edge of the infinite cluster, then given that $X_0,\ldots,X_{i-1}$ is already defined, let $X_i$ be the lexicographically smallest one among all vertices $v$ of the remainder of the infinite cluster of $o$ such that $(X_{i-1},v)$ is an edge of the $1$-UnG and there is a path in the $1$-UnG from $v$ to infinity that does not use any of the edges $(X_0,X_1),\ldots,(X_{i-1},v)$. (Note that there may be several $v$ satisfying this property in case there are multiple edge-disjoint incoming paths of the $1$-DnG at $X_{i-1}$.) If $o$ is adjacent to a bidirectional edge, then let $X_1$ be its neighbor via this edge, and define $X_2,X_3,\ldots$ as before.

Now, if $X_0=o,X_1,X_2$ is the sequence of vertices in an infinite path (on the event $A$), we put
\[ K=\inf \{ k \geq 0 \colon (X_k,X_{k+1}) \text{ is not an edge of the $1$-DnG} \}. \]
We know from (the proof of) Lemma~\ref{lem:surenoperc} that $K<\infty$ a.s. Next, let us define
\[ L=\inf \{ k > K \colon 
(X_{k},X_{k+1})\text{ is an edge of the 1-DnG}\}. \]
Now we claim that $L=\infty$ on the event $\{ K<\infty\}$ (where we put $K=\infty$ on the event $A^c$). Indeed, for all $k=K,K+1,\ldots,L-1$, $(X_{k+1},X_k)$ is an edge of the $1$-DnG because it follows from the definition of $K$ that $(X_{k},X_{k+1})$ is not an edge of the 1-DnG, but $(X_k,X_{k+1})$ is an edge of the 1-UnG, and this is only possible if $(X_{k+1},X_k)$ is an edge of the $1$-DnG. But now, if $L<\infty$, then $(X_{L},X_{L-1})$ is an edge of the $1$-DnG, but then $(X_{L},X_{L+1})$ cannot be an edge of the $1$-DnG since there is only one edge going out of $X_L$. This implies the claim.

\color{blue} Note that this implies that the cluster of $o$ in the $1$-UnG consists of coalescing (finite or infinite) directed paths towards $X_K$ and $X_{K+1}$ if $K\geq 1$; if $K=0$, then all are directed towards $o$. In particular, there can be at most one bidirectional edge in the infinite cluster, namely the one between $X_{K-1}$ and $X_K$ (because if two vertices have outgoing edges of the $1$-DnG towards each other, they cannot have any other outgoing edges, so that $K$ is at most the larger one among the indices of the two endpoints of the bidirectional edge). We perhaps don't need this, but I wanted to note it. \color{black}

Let us denote the $\ell^1$-sphere of radius $n\in\N_0$ by $S^{(1)}_n$ (with $S^{(1)}_0=\{o\}$). From Lemma~\ref{lem:surenoperc} we even know that $\P(\kappa<\infty\,|\,A)$? where
\[ \kappa = \inf \big\{ k \geq 0 \colon o \not\rightsquigarrow S^{(1)}_k \text{ in the $1$-DnG} \big\}. \]
By continuity of measures, we can find $k_0\in\N$ such that $\P(\kappa \leq k_0 |A) \geq 1/2$. Now, by our previous observation, since from $o$ to $S^{(1)}_\kappa$ there always exists a directed path in the $1$-DnG, 
on the event $\{ \kappa \leq k_0 \}$, for all $n\in\N_0$ there must exist a self-avoiding directed path of length $n$ in the $1$-DnG ending at some vertex of $S^{(1)}_{k_0}$. Now, writing $c_{n,k}(d)$ for the number of self-avoiding paths in the $d$-dimensional lattice of length $n$ starting from (or ending at) $S^{(1)}_{k_0}$, it is clear that
\[ \lim_{n\to\infty} c_{n,k}(d)^{1/n} = \lim_{n\to\infty} c_{n,0}(d)^{1/n}=c(d) \leq 2d-1, \]
where $c(d)=\lim_{n\to\infty} c_{n,0}(d)^{1/n}$ is the connective constant of $\Z^d$. Now, any self-avoiding path of length $n$ ending at $S^{(1)}_{k_0}$ is open with probability $(2d)^{-n}$. Hence,
\[ \P(\exists\text{ a self-avoiding path of length $n$ starting from $S^{(1)}_{k_0}$ included in the $1$-DnG}) \leq \Big( \frac{2d-1}{2d} \Big)^n, \]
which decays exponentially fast in $n$. This contradicts the assertion that the left-hand side is bounded from below by $\P(A)/2>0$ independently of $n$. Therefore, $\P(A)=0$, as wanted.
\end{proof}

\subsection{Results for the bidirectional $k$-neighbor graph}
Let us denote by $c(d)$ the connective constant of $\Z^d$, that is $$c(d)=\lim_{n\uparrow\infty}c_n(d)^{1/n},$$
where $c_n(d)$ is the number of self-avoiding paths in the $d$-dimensional cubic lattice, starting at the origin, of length $n$. The limit exists by subadditivity. Also, $d<c(d)<2d-1$.  
\begin{lemma}
For all $k$ such that $k(k-1)<2d(2d-1)/c(d)$ we have that $\theta^{\rm B}(k,d)=0$.
\end{lemma}
Let us note that this results can be interpreted in two ways. First, for given $d$, absence of percolation is guaranteed for sufficiently small $k$, and for example $\theta^{\rm B}(2,d)=0$ for any $d\geq 2$ (while it is clear that $\theta^{\rm B}(2,1)=1$). However, also for fixed $k$, since $2d(2d-1)/c(d)>2d$, for sufficiently large $d$, there is no percolation. This is due to the fact that in high dimensions, it is unlikely for two neighboring vertices to pick the connecting edge. 
\begin{proof}
Note first that the system is negatively correlated in the sense that for all directed edges $\ell_1=(x_1,x_2)\neq\ell_2=(y_1,y_2)$ we have that
$$
\mathbb P(\ell_1\text{ open and }\ell_2\text{ open})\le \mathbb P(\ell_1\text{ open})^2.
$$
Indeed, note that if $x_1\neq y_1$, then the inequality is an equality by independence. However, if $x_1=y_1$ then 
$$
\mathbb P(\ell_1\text{ open}|\ell_2\text{ open})=\frac{k-1}{2d-1}\le \frac{k}{2d}=\mathbb P(\ell_1\text{ open}).
$$
Using the first moment method, we have that 
\begin{align*}
\P(o\leftrightsquigarrow_{\rm B}\partial B_n)\le \sum_{s\in \mathcal R_n}\P(s\text{ is open})\le C (c(d)p)^n 
\end{align*}
where $p=k(k-1)/(2d(2d-1))$ is the probability that the origin chooses precisely two prescribed edges. This leads to the criterion for absence of percolation $c(d)p<1$.
\end{proof}

We need some result on nontrivial percolation behavior
\begin{prop}\label{prop-1indep}
\begin{enumerate}[(i)]
\item For any $d \geq 2$, if $k>2\sqrt{0.8457}d \approx 1.8392d$, then $\theta^{\rm B}(k,d)>0$.
\item For any $\alpha > 2 \sqrt{0.5847} \approx 1.5293$, we have $\theta^{\rm B}(\alpha d, d)>0$ for all $d$ sufficiently large.
\end{enumerate}
\end{prop}
\begin{proof}
Note that the $k$-BnG model (just as the $k$-UnG model) in any dimension $d$ is an 1-independent Bernoulli bond percolation model where each edge is open with probability $p=k^2/(4d^2)$. Let us write 
\[ p_{\max}(\Z^d)= \sup \{ p \in (0,1) \colon \text{some 1-independent Bernoulli bond perc.\ model does not percolate} \}. \]
According to [Balister et al.], we have $p_{\max}(\Z^d) \leq p_{\max}(\Z^{d-1})$ and $p_{\max}(\Z^2) \leq 0.8457$. These assertions imply (i). Part (ii) follows from the assertion of the same paper that $\lim_{d\to\infty} p_{\max}(\Z^d) \leq 0.8457$. 
\end{proof}
\begin{remark}
The smallest nontrivial (i.e., such that $k \neq 2d$) case covered by Proposition~\ref{prop-1indep} is $k=13$ and $d=7$: Since $13/7 >1.8392$, it follows that $\theta^{\rm B}(13,7)>0$. Similarly, $\theta^{\rm B}(24,13)>0$ is the first case where percolation follows for $k=2d-2$, and so on.

Although the $k$-UnG model is also an 1-independent Bernoulli bond percolation model, the same approach unfortunately does not yield any new results in its case. Indeed, here each edge is open with probability $p=k(4d-k)/(4d^2)$. Thus, for $k=d$ it always holds that $p=3/4<0.8457$, while for $k>d$ we already know that $\theta^{\rm D}(k,d)=1$. 
\end{remark}

%\subsection{Results for the exclusively unidirectional $k$-neighbor graph}
%Since the $k$-XnG is a subgraph of the $k$-UnG, the following assertion is immediate.
%\begin{prop}
%We have that $\theta^{\rm X}(1,d)=0$ for all $d \geq 1$. 
%\end{prop}



%Using a similar argument, we can also prove the followin
%\begin{lemma}
%We have $\theta(k,d)\leq \theta (k,d')$ whenever $k\geq 1$ and $d\leq d'$.
%\end{lemma}
%\begin{proof}
%The lemma follows from the fact that the restriction of the $k$-nnG in $d'$ dimensions to a $d$-dimensional coordinate plane is stochastically dominated by a $k$-nnG in $d$ dimensions.
%\end{proof}

\newpage

%%OLD COX-DURRETT BOUND FOR UNDIRECTED/BIDIRECTIONAL, IT HAS ERRORS
%THIS DOESN'T SEEM TO WORK OUT DUE TO PROBLEMS WITH THE MARTINGALES.
%
%With methods analogous to the case of the $k$-UnG, we can additionally prove the following.
%\begin{prop}\label{prop-U}
 %   We have $\theta^{\rm U}(3,4)>0$. Moreover, $\theta^{\rm U}(2,d)>0$ for all $d \geq 5$. \color{red} Need to check issues with lack of independence. \color{black}
%\end{prop}
%\begin{proof}
%We argue analogously to the proof of Theorem~\ref{thm-CoxDurrett}, the only difference being that now the probability of an edge being present in the graph is higher. \color{red} (Add negativity of correlations from the old stuff. \color{black} Namely, we have for an undirected edge $\ell$ with endpoints $x_1,x_2 \in \Z^d$ being nearest neighbors:
%\begin{equation}\label{Uprob}
%\begin{aligned}
%\P&(\ell\text{ is an edge of the $k$-UnG}) = \P((x_1,x_2)\text{ is an edge of the $k$-UnG}) \\ & \quad +\P((x_2,x_1)\text{ is an edge of the $k$-UnG})- \P(\text{both } (x_1,x_2) \text{ and } (x_2,x_1) \text{ are edges of the $k$-UnG} ) \\ &= \frac{k}{2d} + \frac{k}{2d} - \frac{k^2}{4d^2} = \frac{k}{d} \Big( 1 - \frac{k}{4d} \Big),
%\end{aligned}
%\end{equation}
%since the events $\{ (x_1,x_2)\text{ is an edge of the $k$-UnG}\}$ and $\{ (x_2,x_1) \text{ is an edge of the $k$-UnG} \}$ are independent and their probabilities equal $k/(2d)$. Now, using the martingale $W_n=N_n/((\frac{k}{d} ( 1 - \frac{k}{4d} ))^n)$, if $\varrho(d)< \frac{k}{d} ( 1 - \frac{k}{4d} )$, it follows analogously to the proof of Theorem~\ref{thm-CoxDurrett} that $\theta^{\rm U}(k,d)>0$. 
%
%Let us first consider the case $(k,d)=(3,4)$. In the proof of Theorem~\ref{thm-CoxDurrett} we computed that $\varrho(4)=\P(\tau_4<\infty) < 0.5$, and thus
%\[ \varrho(4) < \frac{k}{d} \Big( 1 - \frac{k}{4d} \Big)= \frac{3}{4} \Big( 1- \frac{3}{4 \cdot 4} \Big) = \frac{39}{64} = 0.609375. \] 
%Thus, $\theta^{\rm U}(3,4)>0$.
%
%Let us now focus on the case $(k,d)=(2,5)$. We also computed in the aforementioned proof that $\varrho(5)<0.28$, and thus
%\[ \varrho(5) < \frac{k}{d} \Big( 1 - \frac{k}{4d} \Big) = \frac{2}{5} \Big( 1- \frac{2}{20} \Big) = 0.36. \]
%This implies that $\theta^{\rm U}(2,5)>0$. For $d=6,7,8,9,10$ it in fact holds that the right-hand side of~\eqref{CDUB} is smaller than $\frac{2}{d} (1-2/(4d))$, and for $d \geq 10$ it is again elementary to show that this right-hand side of~\eqref{CDUB} minus $\frac{2}{d} (1-2/(4d))$ stays negative (given that it is negative for $d=9$). \color{red} (Indeed, I think that we can again argue that the rhs of~\eqref{CDUB} for dimension $d+1$ is at most $d/(d+1)$ times the same for dimension $d$, while $\frac{2}{d+1}(1-2/(4d+1))>\frac{d}{d+1} \frac{2}{d}(1-2/(4d))$. \color{black}) Therefore, $\theta^{\rm U}(2,k)>0$ holds for all $k\geq 5$, as asserted.
%\end{proof}
%
%\color{green} If this is fine, the conjecture is proven apart from $(k,d)=(2,2),(2,3),(3,3),(2,4)$. For $(2,4)$ it is not impossible that a Cox--Durrett style bound can work, but we would need $\P(\tau_4=k)$ at least for all $k\leq 8$. So the problem is similar as in the directed case for $k=3$ and $d=4$. \color{black}
%
%
%\begin{prop}\label{prop-B}
%For $d \geq 5$ and $k \geq \sqrt{6 d}$ and for $d \geq 4$ and $k\geq \sqrt{8d}$ we have $\theta^{\rm B}(k,d)>0$.
%\end{prop}
%This implies e.g.\ that $\theta^{\rm B}(5,5)>0$ and $\theta^{\rm B}(6,4)>0$. \color{green} So far, $k=5$ is the smallest $k$ for which we get percolation for some $d$. \color{black}
%\begin{proof}[Proof of Proposition~\ref{prop-B}]
%We again adapt the proof of Theorem~\ref{thm-CoxDurrett}, similarly to how we did it in the proof of Proposition~\ref{prop-U}. \color{red} Also here we have negative correlations -- need to copy it from old stuff. \color{black} For an undirected edge $\ell$ with endpoints $x_1,x_2 \in \Z^d$ being nearest neighbors we have (cf.~\eqref{Uprob})
%\[ \P(\ell \text{ is an edge of the $k$-BnG}) = \P(\text{both } (x_1,x_2) \text{ and } (x_2,x_1) \text{ are edges of the $k$-UnG} )=\frac{k^2}{4d^2}. \]
%Now, we have seen in the aforementioned proofs that $\varrho(d)<\frac{3}{2d}$ for $d \geq 5$ and $\varrho(d)<\frac{2}{d}$ for $d \geq 4$. Consequently, using the martingale $W_n=(4d^2/k^2)^n N_n$, we obtain percolation for $d\geq 5$ whenever 
%\[ \frac{3}{2d} < \frac{k^2}{4d^2}, \]
%which is equivalent to $k \geq \sqrt{6 d}$. For $d \geq 4$ we obtain percolation whenever
%\[ \frac{2}{d} < \frac{k^2}{4d^2}, \]
%in other words, if $k \geq \sqrt{8d}$, as claimed.
%\end{proof}


\section{Percolation in U-NN for $k=d=2$ via dual approach}

\begin{theorem}
    Consider the $2-\mathrm{UnG}$ model in two dimensions. Then, $\theta^U(2,2)>0$. 
\end{theorem}

\begin{proof}
    We prove the positive percolation probability via a dual approach. That is, let $\D^2:=\{(x+\frac12,\frac12): x\in \Z^2\}$ be the two-dimensional dual lattice of $\Z^2$. An edge $e$ in the dual lattice has exactly one edge $e'$ from the lattice $\Z^2$ that crosses $e$, and we say the edge $e$ is closed if and only if the edge $e'$ is open. For the $2-UnG$ model in two dimensions, it thus follows that for any such edge $e$, 
    \be 
    \P(e \text{ is closed})=1-\P(e'\text{ is closed})=1-\frac34=\frac14.
    \ee 
    Moreover, the negative correlations between edge that are present in the $2$-UnG model also appear in its dual. Namely, for two edges in the dual that are vertex disjoint or not orthogonal, the status of the edges is independent. When two edges $e_1,e_2$ in the dual share a vertex and are orthogonal, that is, when there is some vertex $v\in \Z^2$ such that the edges emanating from $v$ can cross both $e_1$ and $e_2$, then
    \be 
    \P(e_1\text{ and }e_2\text{ closed in the dual})=\frac{1}{24}\leq \Big(\frac14\Big)^2=\P(e_1\text{ closed in the dual})\P(e_2\text{ closed in the dual}).
    \ee 
    This can be extended, so that for any $k\geq 2$ and any path $(v_1,\ldots, v_k)$ of nearest-neighbour vertices in $\D^2$, 
    \be 
    \P(\text{Path }(v_1,\ldots, v_k)\text{ closed})\leq \prod_{i=1}^{k-1}\P{(v_i,v_{i+1})\text{ closed}}=\Big(\frac14\Big)^{k-1}.
    \ee 
    As such, we can use the negative correlation to bound the probability that there exists a closed circuit in the dual that surrounds the origin. For any $n\geq 4$, the number of such circuits of length $n$ is at most $n\sigma(n-1)$, where $\sigma(n-1)$ denotes the number of self-avoiding paths in $\Z^2$. As $\sigma(n-1)=(c(d)+o(1))^{n-1}$, we thus arrive at 
    \be \ba\label{eq:finnorpercprob}
    1-\theta^U(2,2)&=\P(\exists \text{ circuit }\gamma \text{ in }\D^2\text{ that surrounds the origin})\\
    &\leq \sum_{n=4}^\infty \sum_{\substack{\text{circuits }\gamma\\ |\gamma|=n}}\P(\gamma\text{ closed in }\D^2)\\
    &\leq \sum_{n=4}^\infty n(c(2)+o(1))^{n-1} \Big(\frac14\Big)^n, 
    \ea \ee 
    which is finite since $c(2)\leq 3$. We now apply an argument from {\color{red} Cite Durret 88 or Jan Swart's notes} to show that this probability is in fact smaller than one, which yields $\theta^U(2,2)>0$. Let $D_m=\{0,\ldots, m\}^2$ and say that a vertex $i\in \Z^2$ is `wet' when there exists a $j\in D_m$ such that $i\to j$, that is, when there exists a path of closed edges from $i$ to $j$. Suppose the component of the origin is finite almost surely. Then, the number of wet vertices is finite almost surely as well, and there exists a circuit in the dual that surrounds all the wet sites. By the same argument as in~\eqref{eq:finnorpercprob}, 
    \be 
    \P(D_m \not\to \infty )\leq \sum_{n=4m}^\infty \sum_{\substack{\text{circuits }\gamma \\ |\gamma|=n}}\P(\gamma\text{ closed in }\D^2)\leq \sum_{n=4m}^\infty n(c(d)+o(1))\Big(\frac14\Big)^n,
    \ee 
    where we can start the outer sum from $4m$ since a circuit that surrounds all wet vertices must surround $D_m$ and hence have length $4m$ at least.\basc{Do we need to allow only vertices in $\Z^2\backslash D_m$ to be wet? Otherwise, if vertices in $D_m$ can be wet (and in particular, it can then be possible that \emph{only} vertices in $D_m$ are wet). If that would be the case, I don't see why such a circuit must surround $D_m$.} 
    As a result, since the sum is finite for each $m\in\N$, choosing $m$ large enough yields that $\P(D_m\to \infty)>0$ and hence $\P(i\to \infty)>0$ for some $i\in D_m$. By translation invariance, it then follows that $\theta^U(2,2)>0$, as desired.
\end{proof}


\newpage
NOT USEFUL, AT LEAST IN THIS CONTEXT:

\medskip
However, when we consider the undirected model in which the directed edges are replaced by undirected edges, we still observe negative correlations, 
$$
\mathbb P(\ell_1\text{ open}|\ell_2\text{ open})=\frac{k-1}{2d-1}+(1-\frac{k-1}{2d-1})\frac{k}{2d}\le \frac{k}{2d}(2-\frac{k}{2d})=\mathbb P(\ell_1\text{ open}).
$$
Similarly, for the bidirectional case
$$
\mathbb P(\ell_1\text{ open}|\ell_2\text{ open})=\frac{k-1}{2d-1}\frac{k}{2d}\le\frac{k^2}{4d^2}=\mathbb P(\ell_1\text{ open}).
$$
In this case, the first moment methods leads to the criterion
$$
k<\frac{2d}{\sqrt{c(d)}}
$$
and thus for example for $k=2$ the process is subcritical for any $d\ge 2$. For example for $k=3$ the process is subcritical for $d\ge 4$, or for $k=4$ the process is subcritical for $d\ge 8$.

Let us note that there is a major qualitative difference between the bidirection graph and the both other graphs. For the bidirectional graph it becomes harder to percolate in high dimensions since neighboring vertices find it harder to pick each other. This seems to be the dominate effect that overcomes the fact that also in higher dimensions it is harder to backtrack or do loops. On the other hand, without the need to make matching picks, the other graphs should percolate easier in higher dimensions since there only the backtracking and looping becomes less likely. 

In fact, we can be better since, in the bidirectional case, if the edges $\ell_1,\ldots,\ell_m$ form a path for $m \geq 2$, we have
$$
\mathbb P(\ell_1\text{ open}|(\ell_i)_{i=2}^{m-1} \text{ open})\le \mathbb P(\ell_1\text{ open}|\ell_2\text{ open})=\frac{k-1}{2d-1}\frac{k}{2d}
$$
and hence the criterion is
$$
\frac{k-1}{2d-1}\frac{k}{2d}c(d)<1.
$$
Since $c(d)<2d-1$, we have $k(k-1)\le 2d$, we have that for $k=3$ and $d=3$, the system is also subcritical. Hence we have the following theorem.
\begin{theorem}
Consider the bidirectional model with $k=3$. Then for $d=2$ the model is supercritical and for $d\ge 3$ the model is subcritical. 
\end{theorem}

Let us note that there are some fundamental differences between the lattice model and the continuum model in the bidirectional case. This can be seen for example by considering the expected degree of the origin. In the discrete case we have
$$
\mathbb E[N_k]=k\frac{k}{2d}, 
$$
which can be seen as the expected value of a Bernoulli experiment in which for each of the $k$ outgoing edges of the origin, we flip a coin to see if this edge is also an ingoing edge. In particular, this expected value tends to zero as $d$ tends to infinity. On the other hand in the continuum by the FKG inequality,
\begin{align*}
\mathbb E[N_k]&=\int_{\mathbb R^d}d x\mathbb P(X(B_{|x|}(o))\le k-1, X(B_{|x|}(x))\le k-1)\\
&\ge \int_{\mathbb R^d}d x\mathbb P(X(B_{|x|}(o))\le k-1)^2\\
&= \int_{\mathbb R^d}d x {\rm Pois}_{v_d|x|^d}(\{0,\dots,k-1\})^2\\
&= \int_0^\infty d r {\rm Pois}_{r}(\{0,\dots,k-1\})^2\\
&= \sum_{i,j=0}^{k-1}\frac{(i+j)!}{i!j!}2^{-1-i-j}
\end{align*}
where we also used the representation of the gamma function. This expression is independent of the dimension and for $k=2$ we have $\mathbb E[N_k]\ge \color{red} 5/4 \color{black}$. Since $\mathbb E[N_k]$ is increasing in $k$, we can not use the dimension to reduce the expected degree. Actually, in the limit $d$ to infinity, the FKG inequality becomes an equality. Also, from simulations, it should be true that $\lim_{k\uparrow\infty}\mathbb E[N_k]/k=1$, so in high dimensions the bidirectional model behaves more like the directional model where  $\mathbb P(N_k=k)=1$. It is also immediate that for the undirected case we have $\mathbb E[N_k]\le 2k$. What is the role of dimension in the continuum model?

For the record, the computation can be made precise without FKG, for the bidirectional model,
\begin{align*}
\mathbb E[N_k]&=\sum_{i=0}^{k-1}\sum_{j=0}^{k-1-i}\sum_{l=0}^{k-1-i}\frac{(i+j+l)!}{i!j!l!}\frac{\hat v_d^i\bar v_d^{j+l}}{(\hat v_d+2\bar v_d)^{1+i+j+l}}
\end{align*}
where $\hat v_d=|B^d_1(o)\cap B^d_1(e_1)|/|B^d_1(o)|$ and $\bar v_d=|B^d_1(o)\setminus B^d_1(e_1)|/|B^d_1(o)|$.

Let $u\sim v$ denote that two vertices $u,v\in \Z^d$ are nearest-neighbours and let $v\to u$ denote that $v$ has an outgoing edge directed towards $u$ in $k$-nn$G$ on $\Z^d$. 

\begin{lemma}[Expected number of neighbours further away from the origin]
	Fix $k\leq d$ and let $v\in \Z^d$ be a vertex such that $|\{u\in \Z^d: u\sim v\text{ and } \|u\|_1=\|v\|_1+1\}|=d$. Then, 
	\be 
	\E{|\{u\in \Z^d: v\rightarrow u\text{ and } \|u\|_1=\|v\|_1+1\}|}=\frac k2.
	\ee 
\end{lemma}

\begin{proof}
	As $v$ has exactly $d$ many nearest-neighbours such that $\|u\|_1=\|v\|_1+1$, we can write the expected value as 
	\be
	\sum_{j=0}^k j\frac{\binom{d}{j}\binom{d}{k-j}}{\binom{2d}{k}}=\frac{d}{\binom{2d}{k}}\sum_{j=0}^k \binom{d-1}{j-1}\binom{d}{k-j}=\frac{d}{\binom{2d}{k}}\sum_{j=0}^{k-1}\binom{d-1}{j}\binom{d}{(k-1)-j}.
	\ee 
	We then use the Vandermonde identity to arrive at 
	\be 
	\E{|\{u\in \Z^d: v\rightarrow u\text{ and } \|u\|_1=\|v\|_1+1\}|}=\frac{d\binom{2d-1}{k-1}}{\binom{2d}{k}}=\frac k2,
	\ee 
	as desired.
\end{proof} 

\subsection{Some ideas for a proof} \color{red} Is this useful or can it go away? -- András \color{black}

We consider the growths process starting at the origin $X^0=\{o\}$, $X^1=\{X_1,X_2\}$, where $X_1, X_2$ are the two neighbors of the origin, the two (out of the $4$) neighbors such that there is a directed edge from the origin to that neighbor. In the same vain we define $X^1=\{X_{1,1},X_{1,2},X_{2,1},X_{2,2}\}$ and  $X^n=\{X_{i_1,\dots,i_n}\}$ and note that in this construction we track the path from the origin to that point via the indices. However these paths are not self-avoiding and the set $X^n$ may contain the same point multiple times. Therefor we also consider the sets of newly discovered vertices
$$
Y^n=X^n\setminus X^{n-1} 
$$
and note that of course also $Y^n$ may contain points with different indices that are actual the same. Note that by construction, if the cluster of the origin is finite, then there must exists a generation $n$ with $Y^n=\emptyset$. 

Our starting point is the following standard ansatz
\begin{align*}
1-\theta(2,2)&= \mathbb P(\text{there exists $n$ such that $Y^n=\emptyset$})\\
&\le  \sum_{n\ge 2}\mathbb P(Y^n=\emptyset)\\
&=  \sum_{n\ge 2}\mathbb E[\mathbb P(Y^n=\emptyset|X^{n-1})]. 
\end{align*}
We note that in order for the process to die out in generation $n$ what is important is the {\em free degrees} of the points in $Y^{n-1}$ that is, for $Y_i\in Y^{n-1}$ we are interested in 
$$
N(Y_i, X^{n-1})=\#\{x\in \mathbb Z^2\colon x\in \partial Y_i\setminus X^{n-1}\},
$$
where $\partial Y_i$ denotes the set of $\mathbb Z^d$-neighbors of $Y_i$. We note that almost surely
\begin{align*}
\mathbb P(Y^n=\emptyset|X^{n-1})&=1\{N(Y_i, X^{n-1})\le 2\text{ for all }Y_i\in Y^{n-1}\}\mathbb P(Y^n=\emptyset|X^{n-1})
\end{align*}
and 
\begin{align*}
\mathbb P(Y^n=\emptyset|X^{n-1})= \Big(\frac{1}{6}\Big)^{\sum_{Y_i\in Y^{n-1}}1\{N(Y_i,X^{n-1})=2\}} \Big(\frac{1}{2}\Big)^{\sum_{Y_i\in Y^{n-1}}1\{N(Y_i,X^{n-1})=1\}}.
\end{align*}
Hence, the evolution of the probability of $Y^n=\emptyset$ is determined by the evolution of 
$$
\mu_{n-1}=\sum_{Y_i\in Y^{n-1}}\delta_{N(Y_i,X^{n-1})}\in \mathcal M(\{0,1,2,3\}),
$$
hence 
\begin{align*}
\mathbb P(Y^n=\emptyset|X^{n-1})&=1^{\mu_{n-1}(0)}\Big(\frac{1}{2}\Big)^{\mu_{n-1}(1)}\Big(\frac{1}{6}\Big)^{\mu_{n-1}(2)}0^{\mu_{n-1}(3)}\\
&=:\exp\Big(\sum_{k=1}^3\mu_{n-1}(k)g(k)\Big)=\exp(\mu_{n-1}(g)) =:f(\mu_{n-1})
\end{align*}
and $f$ is decreasing, i.e., for $\mu\le \nu$ (that is $\mu(k)\le \nu(k)$ for all $k\in \{0,\dots, 3\}$) we have $f(\mu)\ge f(\nu)$. In particular, 
\begin{align*}
1-\theta(2,2)&\le  \sum_{n\ge 1}\mathbb E[{\rm e}^{\mu_n(g)}],
\end{align*}
which should be bounded away from one. 
\color{blue} I was hoping for some kind of martingale argument but did not find one yet. -- Bene\color{black}

\begin{comment}{
\subsection{Undirected variants}
Instead of considering the oriented $k$-nearest neighbor graph itself, we can transform it into an undirected graph in an undirected fashion, i.e., connecting two vertices by an edge whenever there is a directed edge between them in at least one direction, or in a bidirectional fashion, i.e., connecting them if and only if there are directed edges between them in both directions. The following assertion is easy to verify. Let us call the former graph the $\mathbf U$-$k$nnG  and the latter one the $\mathbf B$-$k$nnG, respectively. It is trivial that the $\mathbf B$-$1$nnG does not percolate in any dimensions because its degrees are bounded by 1. Further, for $d=1$ and $k=2$, the graph percolates because it includes all nearest-neighbour edges of $\mathbb Z$.
%
\begin{prop} \color{red} obsolete \color{black}
The $\mathbf B$-2nnG does not percolate for $d\geq 2$.
\end{prop}
\begin{proof}
It suffices to verify the assertion for $d\geq 2$.
Since the $\mathbf B$-$k$nnG has degrees bounded by $k$, it follows that for $k=2$ any infinite path in the graph must be self-avoiding. Now, the probability that a fixed nearest-neighbor edge of $\mathbb Z^d$ is included in the $\mathbf B$-2nnG equals 
\[ \Big(\frac{1}{\binom{2d}{2}}\Big)^2 = \frac{1}{d^2}. \]
%
The number of self-avoided paths of length $n \in \mathbb N$ in $\Z^d$ starting from the origin is bounded by $(2d-1)^n$ (\color{red} cf.\ Grimmett \color{black}). 
Now, for a fixed self-avoiding path of length $2n$ consisting of nearest-neighbor edges of $\Z^d$ starting from $o$, consider the event that  
\begin{itemize}
\item for any $k=1,\ldots,n$, the $2k$-th edge of the path (i.e., the one between the $2k-1$-st and the $2k$-th vertex of the path) is included in the $\mathbf B$-2nnG (this has probability $1/\binom{2d}{2}=\frac{1}{d(2d-1)}$), and
\item for any $k=0,\ldots,n-1$, there is a \emph{directed} edge included in the 2-nnG from the $2k$-th vertex of the path to the $2k+1$-st one (this has probability $(2d-1)/\binom{2d}{2} = \frac{1}{d}$).  
\end{itemize}
Note that since the choices of outgoing 2nnG edges of different vertices are independent, the probability that there exists such a self-avoiding path is at most
\[ \frac{(2d-1)^{2n}}{d^{n} (d(2d-1))^n} = \big(\frac{2d-1}{d^2}\big)^n. \] 
Since $(2d-1)/d^2<1$ for any dimension $d\geq 2$, the proposition follows.
\end{proof}
%
%\bigskip
%Trying to repair our computation from last time, but now let $m$ be the number of corners in a blocking interface of length $l$, then we have $m-4$ anticorners and thus $l-2m+4$ vertices in the interface that are neither corners nor anticorners. Then we have
%\begin{align*}
%\mathbb P(\exists \text{ blocking interface } ) &\leq \sum_{n \geq 2} \mathbb P(\exists \text{ blocking interface through } (n,0) ) \\ &  \leq \sum_{n \geq 2} \sum_{l \geq 2n+2} \sum_{m \geq 1} \mathbb P(\exists \text{ blocking interface through } (n,0) \text{ with} \\ \quad &\text{$m$ corners and $l-2m+4$ vertices at straight lines}) 
%\end{align*}
%\color{red} Sorry, this doesn't seem to give a correct upper bound because requirements on blocking edges are often superfluous, as observed by Bas. -- András \color{black}
%
\section{Proofs}\label{sec:proof}
}
\end{comment}

\newpage 
\section{Branching process attempt \#2}
\textbf{Goal:} Use a branching process as upper-bound to conclude that we actually do not percolate in $d=k=2$. 

\subsection{Notation}
For $n \in \N$ denote by $S_n$ the $\ell^1$-ball in $\Z^2$ and by $\mathcal{R}_n$ the number of directed open paths in $S_n$ that connect the origin to the boundary $\partial S_n$.

\subsection{Heuristics}
If we have a directed open path $\gamma$ in $S_n$ that connects the origin to $\partial S_n$, then usually, its endpoint lies on one of the sides of $\partial S_n$ (and not in one of the four corners). Therefore, its offspring distribution $\mu$ is as follows: 
\begin{align*}
    \mu(0) = \frac{1}{6}, \quad \mu(1) = \frac{2}{3}, \quad \mu(2) = \frac{1}{6}.
\end{align*}
In particular, the typical offspring distribution is critical, and a Branching process with only these typical individuals would die out almost surely. Since it is exactly critical, it will survive for a substantial time which could explain that we don't see this in the simulations. 

Now if we have an open path in $S_n$ that reaches one of the corners of $S_n$, then it will have offspring distribution $\hat{\mu}$ given by 
\begin{align*}
    \mu(0) = 0, \quad \mu(1) = \frac{1}{2}, \quad \mu(2) = \frac{1}{2},
\end{align*}
which is super-critical. 

Now we need to make sure that this branching process dies out almost surely. 







\end{document}

\section{Continuum $k$-nearest-neighbor graphs}
\subsection{Expected degree}
\color{red} What I'm writing here is simple, it may or may not be useful in the end, but I want to write it down before I lose or forget it. -- Bene \color{black}

Let us calculate the expected degree of the typical point in the bidirectional nn-graph based on a Poisson point process
\begin{align*}
\mathbb E\Big[\sum_{X_i\in X}&1\{X(B_{|X_i|}(o))\le k\text{ and }X(B_{|X_i|}(X_i))\le k \}\Big]\\
&=\lambda \int_{\mathbb R^2}{\rm d} x \mathbb P(X(B_{|x|}(o))\le k\text{ and }X(B_{|x|}(x))\le k)\\
&=\lambda \int_0^\infty{\rm d} r 2\pi r \mathbb P(X(B_{r}(o))\le k\text{ and }X(B_{r}(r e_1))\le k)
\end{align*}
Let us note that replacing {\em and} by {\em or} we arrive at a formula for the undirected case. Further, consider the three areas $A_1=B_{r}(o)\setminus B_{r}(re_1)$,  $A_2=B_{r}(r e_1)\setminus B_{r}(o)$ and $B=B_{r}(o)\cap B_{r}(re_1)$ and note that 
$$|B|=r ^2(\frac{\pi}{3}-\frac{\sqrt{3}}{4})=:r^2c$$
and hence $|A_1|=|A_2|=r^2(\pi-c)$. Now let $V_1,V_2$ be two independent Poisson random variables with parameter $\lambda|A_i|$ and $W$ another independent Poisson random variable with parameter $\lambda|B|$. Then, 
\begin{align*}
\mathbb P(X(B_{r}(o))\le k\text{ and }X(B_{r}(r e_1))\le k)=\mathbb P(V_1+W\le k-1\text{ and }V_2+W\le k-1)
\end{align*}
Let us for example calculate the case $k=2$. 
\begin{align*}
\mathbb P(V_1+W\le 1\text{ and }V_2+W\le 1)=\sum_{n=0}^1 \mathbb P(W=n)\mathbb  P(V_1\le 1-n)^2\approx 1.5
\end{align*}

\subsection{Subcriticality for bidirectional nearest-neighbor graphs in high dimensions}
Heuristically, it should be difficult for any pair of points to choose each other as neighbors if the dimension is high. So, let us try a Gouere-type argument. Consider the event $G_\alpha$ that the origin is connected to $B_{10\alpha}\setminus B_{9\alpha}$ only using points in $B_{10\alpha}$. Later it will become important to have that $\mathbb P(G_\alpha)\le 1/2$ for $\alpha\in [1,10]$. Using that 
$$\mathbb P(G_{10\alpha})\le \mathbb P(G_{\alpha})\le 1-\exp(-\lambda|B_{\alpha}|)$$
and the fact that the percolation behavior of our model is independent of $\lambda$, we can choose $\lambda=\log 2/|B_\alpha|$ for a suitable $\alpha$ and this is satisfied.

In order for the argument to work, we need to make the probabilities of following error events small, uniformly in $\alpha$:
\begin{align*}
H_\alpha&=\{\text{ there exists a directly connected pair $x,y\in B_{100\alpha}$ with $|x-y|>\alpha$}\}\\
F_\alpha&=\{\text{ there exists a directly connected pair $x,y$ with $x\in B^c_{10\alpha}$ and $y\in B_{9\alpha}$}\}. 
\end{align*}
Using a first moment bound, we have that 
\begin{align*}
\mathbb P(H_\alpha)&\le \lambda^2\int_{B_{100\alpha}}d x\int_{B_{100\alpha}}d y1\{|x-y|>\alpha\}\mathbb P(x\leftrightarrow y)\\
\mathbb P(F_\alpha)&\le \lambda^2\int_{B^c_{10\alpha}}d x\int_{B_{9\alpha}}d y\mathbb P(x\leftrightarrow y)
\end{align*}
Note that $\mathbb P(x\leftrightarrow y)=f_k(\lambda|B_{|x-y|}|)$ since we are dealing with Poisson probabilities. Using the form of $\lambda$ we hence have that 
\begin{align*}
\mathbb P(H_\alpha)&\le c_1\int_{B_{100}}d x\int_{B_{100}}d y1\{|x-y|>1\}f_k(|B_{|x-y|}|)\\
\mathbb P(F_\alpha)&\le c_2\int_{B^c_{10}}d x\int_{B_{9}}d yf_k(|B_{|x-y|}|)
\end{align*}

\newpage

\section{Some further brainstorming}
From Bene:
\begin{itemize}
\item The B-kNN graph is related to SIR graphs by the fact that in both cases the edge-drawing mechanism depends not only on the relation of two points but also on other points
\item The B-kNN graph based on Poisson points should be stabilizing
\item One could consider other Poisson--Boolean models in which the ball around each point is random and depends on the surrounding point configuration.
\item The B-kNN graph is also somehow related to Poisson (directed) navigations.
\end{itemize}
Some more (from András):
\begin{itemize}
\item We see that even if we knew that there was an infinite cluster containing infinitely many pairwise disjoint infinite paths (or: ``the infinite cluster'' if we also had uniqueness), there must not be an essential trifurcation point, due to possible spider webs. But at least, there exists a degree-3 point with 3 edge-disjoint paths going out from it to infinity, call such a point \emph{good}. So there must be at least one infinite cluster with infinitely many good points. Now, I hope that for each good point we can find a spanning tree of the infinite cluster within which the good point is an essential trifurcation point. If we can even show that in any (nice) spanning tree the essential trifurcation points have a positive density, we can maybe get back to the Burton--Keane style argument that implies that the boundary of the graph in a ball/box is too large. \\
That is, I suspect that the spider-web is just some kind of decoration of an infinite 3-uniform tree.
\item In any case, such a line of argument requires the existence of an infinite cluster with infinitely many pairwise disjoint infinite paths, which I don't know how to prove, but I still hope that looking at half-spaces/sectors is the easiest way to attack this problem.
\end{itemize}

\subsection{A possible line of attack}
Let us try to reprove absence of percolation in the Bk-NN graph for $k=2$ via the following branching method. Let us denote by $B_r(x)$ the open ball with radius $r$ centered at $x$ and by $D_r(x)=B_r(x)\setminus \{x\}$, the ball with removed center. Then, we have for any $\alpha> 0$ we have
\begin{align*}
&\mathbb P(o\leftrightsquigarrow\infty)\le \mathbb P(\text{there exists a self-avoiding path starting at $o$ and ending in $B^{\rm c}_\alpha(o)$})\\
&\le \sum_{n\ge 1}\mathbb E\Big[\sum_{X_1,\dots, X_n\in X}^{\neq}\prod^n_{i=1}1\{X(D_{|X_i-X_{i-1}|}(X_{i-1}))\le 1,X(D_{|X_i-X_{i-1}|}(X_{i}))\le 1\}1\{X_n\in B^{\rm c}_\alpha(o)\}\Big]
\end{align*}
where $X_0\equiv o$ denotes the origin. The bound $\le 1$ indicates that in the open balls with removed centers, there can be at most one additional point. When moving to larger $k$, such as $k=3$, this bound changes accordingly. The idea is that for small $n$, the expectations are small because for large $\alpha$ we require large void spaces. On the other hand, for fixed $\alpha$ and large $n$, in order not to produce loops, the forbidden areas become also large. 

We can use the Slyvniak--Mecke formula to express, for any $n\ge 1$
\begin{align*}
&\mathbb E\Big[\sum_{X_1,\dots, X_n\in X}^{\neq}\prod^n_{i=1}1\{X(D_{|X_i-X_{i-1}|}(X_{i-1}))\le 1,X(D_{|X_i-X_{i-1}|}(X_{i}))\le 1\}1\{X_n\in B^{\rm c}_\alpha(o)\}\Big]\\
&=\int{\rm d} x_1\dots\int{\rm d} x_{n-1} \int_{B^{\rm c}_\alpha(o)}{\rm d} x_n\\
&\hspace{1cm}\mathbb E\Big[\prod^n_{i=1}1\{(X\cup\{x_1,\dots,x_n\})(D_{|x_i-x_{i-1}|}(x_{i-1}))\le 1,(X\cup\{x_1,\dots,x_n\})(D_{|x_i-x_{i-1}|}(x_{i}))\le 1\}\Big].
\end{align*}
Let us note that in the events $\{(X\cup\{x_1,\dots,x_n\})(D_{|x_i-x_{i-1}|}(x_{i-1}))\le 1\}$ on the right-hand side, the additional points $x_{i-1},x_i$ are irrelevant. This method, if successful, can also be used to say something about the decay of the finite-volume percolation function. 

Now it becomes tedious, 
\begin{align*}
&\mathbb E\Big[\prod^n_{i=1}1\{(X\cup\{x_1,\dots,x_n\})(D_{|x_i-x_{i-1}|}(x_{i-1}))\le 1,(X\cup\{x_1,\dots,x_n\})(D_{|x_i-x_{i-1}|}(x_{i}))\le 1\}\Big]\\
&=\mathbb E\Big[\Big(\sum_{A_1\subset [n]}\prod_{i\in A_1}1\{(\{x_1,\dots,x_n\})(D_{|x_i-x_{i-1}|}(x_{i-1}))=1\}1\{X(D_{|x_i-x_{i-1}|}(x_{i-1}))=0\}\times\\
&\qquad\qquad\qquad\prod_{i\in A^c_1}1\{(\{x_1,\dots,x_n\})(D_{|x_i-x_{i-1}|}(x_{i-1}))=0\}1\{X(D_{|x_i-x_{i-1}|}(x_{i-1}))\le1\}\Big)\times\\
&\qquad\Big(\sum_{A_2\subset [n]}\prod_{i\in A_2}1\{(\{x_1,\dots,x_n\})(D_{|x_i-x_{i-1}|}(x_{i}))=1\}1\{X(D_{|x_i-x_{i-1}|}(x_{i}))=0\}\times\\
&\qquad\qquad\qquad\prod_{i\in A^c_2}1\{(\{x_1,\dots,x_n\})(D_{|x_i-x_{i-1}|}(x_{i}))=0\}1\{X(D_{|x_i-x_{i-1}|}(x_{i}))\le1\}\Big)\Big]
\end{align*}
which is equal to 
\begin{align*}
&\sum_{A_1,A_2\subset [n]}\prod_{i\in A_1}1\{(\{x_1,\dots,x_n\})(D_{|x_i-x_{i-1}|}(x_{i-1}))=1\}\prod_{i\in A^c_1}1\{(\{x_1,\dots,x_n\})(D_{|x_i-x_{i-1}|}(x_{i-1}))=0\}\times\\
&\qquad\qquad\prod_{i\in A_2}1\{(\{x_1,\dots,x_n\})(D_{|x_i-x_{i-1}|}(x_{i}))=1\}\prod_{i\in A^c_2}1\{(\{x_1,\dots,x_n\})(D_{|x_i-x_{i-1}|}(x_{i}))=0\}\times\\
&\qquad\mathbb E\Big(\prod_{i\in A_1}1\{X(D_{|x_i-x_{i-1}|}(x_{i-1}))=0\}\prod_{i\in A^c_1}1\{X(D_{|x_i-x_{i-1}|}(x_{i-1}))\le1\}\times\\
&\qquad\qquad\prod_{i\in A_2}1\{X(D_{|x_i-x_{i-1}|}(x_{i}))=0\}\prod_{i\in A^c_2}1\{X(D_{|x_i-x_{i-1}|}(x_{i}))\le1\}\Big)
\end{align*}
So what happens here is that we prescribe deterministic points $x_1,\dots,x_n$ and via $A_1,A_2$ specify that for example $x_i$ is the second nearest neighbor of $x_{i-1}$ and $x_{i-1}$ is the first or second nearest neighbor of $x_{i}$. This sequence of deterministic points has to be legitimate in the sense that the points are in a configuration that is not impossible to achieve by putting additional (random) points. Now, given this sequence of prescribed points, we have to evaluate the probability to see the random points to be positioned precisely such that the sequence is legitimate. Let us note that that points may need to be places in order to make for example some deterministic point the second nearest neighbor. In the end, the probabilities can be evaluated by sums of independent Poisson random variables with parameters given by areas of intersecting balls. 

Next step:
\begin{align*}
&\mathbb E\Big(\prod_{i\in A_1}1\{X(D_{|x_i-x_{i-1}|}(x_{i-1}))=0\}\prod_{i\in A^c_1}1\{X(D_{|x_i-x_{i-1}|}(x_{i-1}))\le1\}\times\\
&\qquad\qquad\prod_{i\in A_2}1\{X(D_{|x_i-x_{i-1}|}(x_{i}))=0\}\prod_{i\in A^c_2}1\{X(D_{|x_i-x_{i-1}|}(x_{i}))\le1\}\Big)\\
&=\sum_{B_1\subset A_1^c}\sum_{B_2\subset A_2^c}\mathbb E\Big(\prod_{i\in A_1\cup B_1}1\{X(D_{|x_i-x_{i-1}|}(x_{i-1}))=0\}\prod_{i\in A^c_1\setminus B_1}1\{X(D_{|x_i-x_{i-1}|}(x_{i-1}))=1\}\times\\
&\qquad\qquad\prod_{i\in A_2\cup B_2}1\{X(D_{|x_i-x_{i-1}|}(x_{i-1}))=0\}\prod_{i\in A^c_2\setminus B_2}1\{X(D_{|x_i-x_{i-1}|}(x_{i-1}))=1\}\Big)\\
&=\sum_{B_1\subset A_1^c}\sum_{B_2\subset A_2^c}\mathbb E\Big(1\{X(\bigcup_{i\in A_1\cup B_1\cup A_2\cup B_2}D_{|x_i-x_{i-1}|}(x_{i-1}))=0\}\times\\
&\qquad\qquad \prod_{i\in (A^c_1\setminus B_1)\cup (A^c_2\setminus B_2)}1\{X(D_{|x_i-x_{i-1}|}(x_{i-1}))=1\}\Big)\\
&=\sum_{B_1\subset A_1^c}\sum_{B_2\subset A_2^c}\exp\Big(-|\bigcup_{i\in A_1\cup B_1\cup A_2\cup B_2}D_{|x_i-x_{i-1}|}(x_{i-1})|\Big)\times\\
&\qquad\qquad \mathbb E\Big[\prod_{i\in (A^c_1\setminus B_1)\cup (A^c_2\setminus B_2)}1\{X(D_{|x_i-x_{i-1}|}(x_{i-1})\setminus \bigcup_{j\in A_1\cup B_1\cup A_2\cup B_2}D_{|x_j-x_{j-1}|}(x_{j-1}))=1\}\Big]\\
\end{align*}
{\color{blue} Bas: Isn't the second line of the second equation incorrect? I think it should be $D_{|x_i-x_{i-1}|}(x_i)$ instead of $D_{|x_i-x_{i-1}|}(x_{i-1})$. And if so, the third and fourth step need correcting, too, I'd say.
}

\color{red}
A bit tighter upper bound, which might be a bit easier to compute because in many areas we just have void spaces after applying the Slivnyak--Mecke formula and thus we don't need much inclusion-exclusion (let me abuse notation and use $X$ both as a set and as the number of Poisson points in the set):
\begin{align*}
&\mathbb P(o\leftrightsquigarrow\infty)\le \mathbb P(\text{there exists a self-avoiding path starting at $o$ and ending in $B^{\rm c}_\alpha(o)$} \\ &\qquad \text{whose vertices do not form a loop in the whole $\mathbf B$-2NN graph} )\\
&\le \sum_{n\ge 1}\mathbb E\Big[\sum_{X_1,\dots, X_n\in X}^{\neq}\prod^n_{i=1}1\{X \cap D_{|X_i-X_{i-1}|}(X_{i-1}) \subseteq \{ X_{i-2} \}, \forall i=2,\ldots,n, \\ & X\cap D_{|X_i-X_{i-1}|}(X_{i}) \subseteq \{ X_{i+1} \}, \forall i=1,\ldots,n-1, \\& X(D_{|X_1-X_{0}|}(X_0)) \leq 1, X_n \notin D_{|X_1-X_{0}|}(X_i) , X(D_{|X_n-X_{n-1}|}(X_n)) \leq 1, X_0 \notin D_{|X_n-X_{n-1}|}(X_n) \} \Big]
\end{align*}
Using the Slivnyak--Mecke formula, for fixed $n$ we can express the expectation as
\begin{align*}
\int& {\rm d} x_1 \ldots \int {\rm d} x_{n-1} \int_{B_\alpha^c(o)} {\rm d} x _n \\ & \mathds E \Big[ \prod_{i=2}^n \mathds 1 \big\{ (X \cup \{ x_1,\ldots,x_n\}) (D_{|x_i-x_{i-1}|}(x_{i-1}) \subseteq \{ x_{i-2} \} \big\} \\
&\times \prod_{i=1}^{n-1} \mathds 1 \big\{ \{ (X \cup \{ x_1,\ldots,x_n\}) (D_{|x_i-x_{i-1}|}(x_{i}) \subseteq \{ x_{i+1} \} \big\} \\&
\qquad \times \mathds 1 \{ x_n \notin (X \cup \{ x_1,\ldots,x_n\}) (D_{|x_1-x_{0}|}(x_{0}), (X \cup \{ x_1,\ldots,x_n\}) (D_{|x_1-x_{0}|}(x_{0}) \leq 1 \} \\ &\qquad \times \mathds 1 \{ o \notin (X \cup \{ x_1,\ldots,x_n\}) (D_{|x_n-x_{n-1}|}(x_{n-1}), (X \cup \{ x_1,\ldots,x_n\}) (D_{|x_n-x_{n-1}|}(x_{n-1}) \leq 1 \} \Big].
\end{align*}

\newpage


\section{General graphs: trifurcation}
\color{red} What I'm writing here is simple, it may or may not be useful in the end, but I want to write it down before I lose or forget it. -- András \color{black}

Let $G=(V,E)$ be a countably infinite connected simple graph with degrees bounded by 3. An \emph{essential trifurcation point} is a vertex $v \in G$ such that the graph with vertex set $V \setminus \{ v \}$ and edge set $E|_{(V \setminus \{ v\})^2}$ consists of three infinite connected components. Of course, essential trifurcation points have degree 3.

\begin{lemma}\label{lem-loop}
No essential trifurcation point of $G$ is included in a loop of the graph.
\end{lemma}
\begin{proof}
Let $v \in V$ be a degree-3 point included in a loop. Let $v_1,v_2,v_3$ be its neighbours. At least two of these, say $v_1,v_2$, are included in the loop. Hence, there exists a path between $v_1$ and $v_2$ not including $v$. But then, after removing $v$ and the edges $(v,v_1),(v,v_2),(v,v_3)$ from the graph, $v_1$ and $v_2$ are still in the same connected component. Hence, $v$ cannot be an essential trifurcation point.
\end{proof}
\begin{cor}\label{cor-almosttree}
Let $v,w \in G$ be two different essential trifurcation points. Then there exists a unique pair $(v_1,w_1) \in G\times G$ such that any path between $v$ and $w$ starts with the edge $(v,v_1)$ and ends with the edge $(w_1,w)$.
\end{cor}
Note that the case when $v_1=w$ and $w_1=v$ is not excluded here.
\begin{proof}[Proof of Corollary~\ref{cor-almosttree}]
Assume otherwise, so that without loss of generality there exist $v_1,v_2$ adjacent to $v$ and one path from $v$ to $w$ starting with the edge $(v,v_1)$ and one starting with the edge $(v,v_2)$. If these two paths are edge-disjoint, then the union of their edges forms a loop including both $v$ and $w$, which contradicts with Lemma~\ref{lem-loop}. Otherwise, the two paths have a joint vertex $z$. Thus, the path consisting of the edges of the path from $v$ to $w$ starting with $(v,v_1)$ up to $z$ and the edges of the path from $z$ to $v$ form a loop including $v$, another contradiction with the lemma.
\end{proof}
Iterating this corollary, we obtain the following assertion.
\begin{cor}
Let $v,w$ be two different essential trifurcation points of $G$. Then, there exists a unique $n\in\mathbb N$ and $(v_0, v_1,\ldots,v_n) \in V^{n+1}$ mit $v_0=v$ and $w_0=w$ such that all the $v_i$ are essential trifurcation points and any path between $v$ and $w$ goes through the points $v_0,v_1,\ldots,v_n$ in this order.
\end{cor}

%Thus, for any two different essential trifurcation points $v,w \in V$ there exists a unique pair of vertices $(s(v,w),t(v,w)) \in G\times G$ such that any path from $v$ to $w$ starts with the edge $(v,s(v,w))$ and ends with the edge $(t(v,w),w)$. Of course, $s(v,w)=t(w,v)$ and $s(w,v)=t(v,w)$. We say that $v,w$ are \emph{almost direct neighbours} if there exists a path from $v$ to $w$ that does not contain any essential trifurcation point apart from $v$ and $w$. 

The \emph{essential tree} $T(G)=(V(T(G)),E(T(G)))$ of $G$ is the graph with vertex set $V(T(G))=\{v \in V \colon v \text{ essential trifurcation point of $G$}\}$ and edge set $E(T(G))=\{ (v,w) \in V(T(G)) \times V(T(G)) \colon v,w\text{ direct neighbours}\}$. 
\begin{lemma}\label{lem-contractedtree}
$T(G)$ is a tree.
\end{lemma}
\begin{proof}
The connectedness of $T(G)$ follows from the fact that for any $v,w \in V(T(G))$, the corresponding path $(v_0, v_1,\ldots,v_n)$ in the statement of Corollary~\ref{cor-almosttree} exists, and the loop-freeness of $T(G)$ follows from the uniqueness of the fact that for any $v,w \in V(T(G))$, the corresponding path is unique.
\end{proof}
\begin{lemma}\label{lem-tree}
Let $H=(V(H),E(H))$ be a tree with degrees bounded by 3, and let $F \subset V(H)$ be a collection of vertices of degree equal to 3. Now, if $|F|=k \in \mathbb N$, then $H$ has at least $2k+2$ vertices. Further, there are at least $\lceil (k+2)/2 \rceil$ vertices in $F$ that are connected by an edge to at least one point in $V(H) \setminus F$.
\end{lemma}
\begin{proof}
Let us call $v,w \in F$ direct neighbours if the (unique) path between them contains no degree-3 point apart from $v$ and $w$. Then, contracting paths between direct neighbours to a single edge, we cannot increase the number of vertices. After all these contractions, the edges between points of $F$ form a tree, which thus has $k-1$ vertices, and each vertex of $F$ still has degree 3. Now, there are $3k$ half-edges adjacent to vertices in $F$, and at most $2k-2$ of them can correspond to edges between vertices of $F$. Hence, there must still be at least $k+2$ half-edges that are adjacent to vertices outside $F$ that have not been contracted. These $k+2$ vertices must be pairwise distinct. Indeed, otherwise some of them , say $z$, would be adjacent to two vertices $v,w$ of $F$, and then $(v,z),(z,w)$ would form a path between $v$ and $z$ in $H$.
This is however impossible because we have another path between $v$ and $w$ in $H$ that only goes through points of $F$ and contracted vertices. We conclude that $H$ must have at least $|F|+k+2=2k+2$ vertices, as wanted.

The same contraction procedure can also not increase the number of vertices in $F$ adjacent to a point in $V(H) \setminus F$. After the contractions, there are precisely $k+2$ vertices outside $F$ adjacent to $F$, as we have seen, while each vertex $v$ of $F$ is connected by at least one edge to $F \setminus \{ v \}$, consequently it is connected to at most 2 vertices outside $F$. By the pigeon-hole principle, we conclude that at least $\lceil (k+2)/2 \rceil$ vertices of $F$ have to be adjacent to vertices outside $F$ in the contracted graph, and thus also in $H$.
\end{proof}
Let us now get back to our infinite graph $G$ studied before. For $W \subseteq V(G)$ define the \emph{inner boundary} of $W$ as
\[ \partial^{\text{in}}~W = \{ w \in W \colon \exists v \in V(G) \setminus W \colon (w,v) \in E(G) \} \]
and
\[ \partial^{\text{out}}~ W = \{ w \in V(G) \setminus W \colon \exists V\in W \colon (w,v) \in E(G) \}. \]
\begin{cor}
Let $W \subset V(G)$ be a set containing only essential trifurcation points of $G$ with $|V| \geq k$. Then, we have
\begin{enumerate}
\item $|\partial^{\text{in}}~W| \geq \lceil (k+2)/2 \rceil$,
\item $|\partial^{\text{out}}~ W| \geq k+2$.
\end{enumerate}
\end{cor}
\begin{proof}
The corollary is a direct consequence of Lemmas~\ref{lem-contractedtree} and \ref{lem-tree}. The tree $F$ corresponding to Lemma~\ref{lem-tree} can be chosen as any finite subtree of $T(G)$ whose vertex set contains $W$.
\end{proof}
\color{red} This should eventually yield us a ``linearly growing boundary'', although I'm not sure yet what to do with those points that are not essential trifurcation points. \color{black}


%\newpage
%\color{black}
%\subsection{Trying to repair the trifurcation argument}
%
%Back to the B-3NN graph, assume that there is an infinite cluster containing infinitely many edge-disjoint infinite paths (this would be true if we knew that infinite clusters ``grow in all directions''). Now, it is easy to show that this cluster contains a \emph{good} point, i.e., a point that has 3 edge-disjoint \color{blue} infinite \color{black} paths coming out of it. Consequently, good points must have a positive density (on the event of percolation if it is not a 0-1 event etc.). Now, we inductively grow a sequence of trees $(T_n)_{n\in\mathbb N_0}$ with infinitely many essential trifurcation points that is a subgraph of the original infinite cluster.
%
%Consider a good point $w_0$ and the three edge-disjoint infinite paths $p_1,p_2,p_3$ going out of it. Initially, $T_0$ consists of the vertices and edges of the union of these three paths. Since the infinite cluster is connected and contains infinitely many edge-disjoint infinite paths, there must exist a good point $w$ in one of the three paths (apart from $w_0$), say in $p_1$. Two edges adjacent to $w$ belong to $p_1$, let $e$ denote the third edge. The question is now whether there exists an infinite path adjacent to $w$ that is edge-disjoint from $p_2$ and $p_3$. Note that if there exists one, it also has to be vertex-disjoint from $p_2$ and $p_3$ because of the degree bound of 3 (and because $w$ is not equal to $w_0$). So if there exists none, this is because all paths starting from $w$ with $e$ share at least one edge with $p_2$ or $p_3$. Otherwise, from $w$ there is a path to infinity starting with $e$ that is edge-disjoint from $p_1,p_2,p_3$. 
%
%Since the cluster contains infinitely many edge-disjoint infinite paths, we conclude that on $p_1,p_2$ or $p_3$ there exists a good point from which there is a path to infinity that is edge-disjoint from $p_1,p_2,p_3$; choose such a point $w_1$ and add the edges and remaining vertices of the infinite path, say $p_4$, to $T_0$. The graph $T_1$ we obtain this way is still a tree.
%
%Iterating this procedure, after $n$ steps we can always find a new point $w_n$ in the current tree $T_n$ from which an infinite path goes out to $\infty$, being edge-disjoint from all the previous $n+2$ infinite paths. After $n$ steps, we have $n$ edge-disjoint infinite paths in $T_n$, and $T_n$ is a tree. So we can take the union $T$ of all these trees. By construction we obtain a tree \color{red} clear? Does connectedness need a proof? Acyclicity needs no proof \color{black}, and the points $w_n$ are always essential trifurcation points of this tree because if one considers two points $x,y$ on two different edge-disjoint paths starting from $w_n$, then for $k\geq n$ the only path in $T_k$ from $x$ to $y$ is via $w_n$, so that if we remove $w_n$, $T_k$ splits into three infinite clusters, and thus so does $T$. 
%
%\color{red} Now, if this really work, we just have to show that conditional on the existence of an infinite cluster, there exists an infinite cluster containing infinitely many infinite paths. The rest should follow from the Burton--Keane argument... \color{black}
\end{document}



%%%%%%%%%%%%%%%%% References %%%%%%%%%%%%%%%%%%
\footnotesize{\printbibliography}
\end{document}
